\section{Domain Task 9: Phân tích danh mục chủ nhà (Host Portfolio Analysis)}

Câu hỏi nghiệp vụ: Nhóm chủ nhà cá nhân và chủ nhà chuyên nghiệp có cơ
cấu phân bố như thế nào, và có sự chênh lệch ra sao về hiệu suất kinh
doanh (điểm đánh giá, tỷ lệ lấp đầy)?

Mục đích: Đánh giá mức độ thương mại hóa của thị trường Airbnb NYC và so
sánh lợi thế cạnh tranh giữa mô hình kinh doanh cá thể so với tổ
chức/chuỗi phòng.

\subsection{Biểu đồ 1: Tỉ trọng cơ cấu loại chủ nhà (Pie Chart)}

\subsubsection*{A. Thiết kế Idiom}

\begin{longtable}{|p{2.8cm}|p{11.5cm}|}
\hline
\textbf{Đặc điểm} & \textbf{Chi tiết} \\
\hline
\endhead
\hline
\endfoot
Idiom & Pie Chart (Biểu đồ tròn) \\
\hline
What & Loại chủ nhà (Host Type): C, Tỷ trọng đóng góp (\% of count): Q \\
\hline
Why & produce (derive) $\rightarrow$ explore $\rightarrow$ summarize \newline
\newline
Derive: phân loại chủ nhà dựa trên số lượng listing sở hữu (1 phòng = Cá nhân, $>$ 1 phòng = Chuyên nghiệp) \newline
\newline
Explore và Summarize: xem nhóm nào đang chiếm ưu thế kiểm soát nguồn cung trên thị trường \\
\hline
How & \textbf{Encode:} \newline
\newline
- Mark: Area (vùng góc quạt) \newline
\newline
- Channel: Angle: Độ lớn của tỷ trọng phần trăm (0\%--100\%) \newline
\newline
\hspace{0.3cm} + Color: Phân biệt 2 nhóm chủ nhà \newline
\newline
\textbf{Manipulate:} \newline
\newline
- Hover + Tooltip: Rê chuột vào từng phần quạt để xem chi tiết tên nhóm và tỷ lệ \% chính xác \newline
\newline
\textbf{Reduce:} \newline
\newline
- Không áp dụng \\
\hline
Scale & - Main key: 2 (Cá nhân, Chuyên nghiệp) \newline
\newline
- Value range: 0\%--100\% \\
\hline
\end{longtable}

\begin{figure}[H]
    \centering
    \safeincludegraphics[width=0.9\linewidth]{images/domain_task9_chart1.png}
    \caption{Hình 9.1. Biểu đồ tỉ trọng cơ cấu loại chủ nhà}
    \label{fig:domain-task9-chart1}
\end{figure}

\subsubsection*{B. Đánh giá biểu đồ}

\textbf{1. Nguyên lý biểu đạt (Expressiveness):}

\begin{itemize}
    \item Thuộc tính: Loại chủ nhà (C), Tỷ trọng (Q).
    \item Channel:
    \begin{itemize}
        \item Angle: Thể hiện tỉ lệ phần trăm $\rightarrow$ dùng cho thuộc tính Q mang tính part-to-whole.
        \item Color (Hue): Phân loại chủ nhà $\rightarrow$ dùng cho thuộc tính Categorical (C).
    \end{itemize}
    \item Đánh giá mức độ quan trọng:
    \begin{itemize}
        \item Ưu tiên 1 -- Góc (Angle): Trong biểu đồ tròn, thuộc tính định lượng (Tỷ trọng \%) được ưu tiên hàng đầu để thấy được độ lớn thị phần. Theo Mackinlay, kênh Angle nằm ở mức khá cho dữ liệu (Q) và là chuẩn mực cho cấu trúc part-to-whole (thành phần trên tổng thể).
        \item Ưu tiên 2 -- Màu sắc (Hue Color): Thuộc tính phân loại (Loại chủ nhà -- C) được gán vào kênh Hue. Theo Mackinlay, Hue là kênh mạnh nhất để biểu diễn dữ liệu Categorical, giúp mắt người phân định ngay lập tức 2 nhóm mà không cần đọc text.
    \end{itemize}
\end{itemize}

\textbf{2. Nguyên lý hiệu quả (Effectiveness):}

\begin{itemize}
    \item Độ chính xác (Accuracy): Kênh Area/Angle có độ chính xác không cao. Việc ước lượng định lượng thông qua Góc/Diện tích có $n \approx 0.7$, độ lỗi Log\_error khá lớn. Mắt người khó so sánh chuẩn xác nếu 2 phần quạt có kích thường gần bằng nhau. Tuy nhiên, ở đây mức chênh lệch khá rõ (62.69\% và 37.31\%), kết hợp thao tác Manipulate (Hover xem Tooltip) đã bù đắp hoàn toàn nhược điểm này.
    \item Khả năng phân biệt (Discriminability): Biểu đồ sử dụng 2 màu tương phản (Xanh/Cam). Số lượng phần loại $< 7$ nên nhận thức màu hoàn toàn không bị ảnh hưởng, rất dễ phân biệt bằng mắt thường.
    \item Khả năng tách biệt (Separability): Kênh góc (Angle) và kênh màu sắc (Color) tách biệt hoàn toàn, không gây nhiễu cho nhau trong quá trình cảm nhận.
\end{itemize}

\subsubsection*{C. Phân tích biểu đồ (Insight)}

\begin{itemize}
    \item Thị trường Airbnb NYC mang tính thương mại hóa cao khi các chủ nhà ``Chuyên nghiệp'' (sở hữu $>$1 phòng) áp đảo nguồn cung với 62.69\% thị phần.
    \item Chủ nhà ``Cá nhân'' (Kinh doanh phòng lẻ) chỉ chiếm phần nhỏ với 37.31\%.
\end{itemize}

\subsection{Biểu đồ 2: So sánh hiệu suất kinh doanh theo loại chủ nhà (Bar Chart)}

\subsubsection*{A. Idiom}

\begin{longtable}{|p{2.8cm}|p{11.5cm}|}
\hline
\textbf{Đặc điểm} & \textbf{Chi tiết} \\
\hline
\endhead
\hline
\endfoot
Idiom & Bar Chart \\
\hline
What & Loại chủ nhà (Host Type): C, Điểm đánh giá (Avg Rating): Q, Tỷ lệ lấp đầy (Avg Is Booked): Q \\
\hline
Why & produce (derive) $\rightarrow$ lookup $\rightarrow$ compare \newline
\newline
derive: tính trung bình cho Rating và tỉ lệ lấp đầy \newline
\newline
Compare: trực diện chất lượng dịch vụ và khả năng thu hút khách hàng giữa 2 nhóm chủ nhà \\
\hline
How & \textbf{Encode:} \newline
\newline
- Mark: Line (Bar mark) \newline
\newline
- Channel: \newline
\newline
\hspace{0.3cm} - PosX: Phân biệt 2 loại chủ nhà \newline
\newline
\hspace{0.3cm} - PosY: thể hiện độ lớn của Avg Rating / Avg Is Booked \newline
\newline
\hspace{0.3cm} - Hue Color: Phân biệt loại chủ nhà \newline
\newline
\textbf{Manipulate:} \newline
\newline
- Hover + Tooltip: Xem chi tiết các con số \newline
\newline
- Facet (Juxtapose): tách thành 2 biểu đồ xếp dọc nhau (1 cho rating, 1 cho IsBooked) dùng chung trục X để dễ đối chiếu \\
\hline
Scale & - Main key: 2 (Loại chủ nhà) \\
\hline
\end{longtable}

\begin{figure}[H]
    \centering
    \safeincludegraphics[width=0.9\linewidth]{images/domain_task9_chart2.png}
    \caption{Hình 9.2. Biểu đồ so sánh hiệu suất kinh doanh theo loại chủ nhà}
    \label{fig:domain-task9-chart2}
\end{figure}

\subsubsection*{B. Đánh giá biểu đồ}

\textbf{1. Nguyên lý biểu đạt:}

\begin{itemize}
    \item Thuộc tính: Loại chủ nhà (C), các chỉ số hiệu suất (Q).
    \item Channel:
    \begin{itemize}
        \item PosX: Định vị nhóm $\rightarrow$ dùng cho thuộc tính (C).
        \item PosY: Chiều dài cột $\rightarrow$ dùng cho thuộc tính định lượng (Q).
        \item Color (Hue): Nhấn mạnh sự phân loại $\rightarrow$ dùng cho thuộc tính (C).
    \end{itemize}
    \item Đánh giá mức độ quan trọng:
    \begin{itemize}
        \item Ưu tiên 1 -- Chiều dài (PosY / Length): Hiệu suất kinh doanh (Q) là thông tin cốt lõi nhất cần so sánh. Theo Mackinlay, Position/Length trên một trục chung là kênh mạnh nhất, chính xác tuyệt đối nhất cho dữ liệu (Q).
        \item Ưu tiên 2 -- Vị trí không gian (PosX): Dùng để chia tách 2 loại chủ nhà (C). Kênh Position cũng là kênh mạnh nhất cho dữ liệu định danh (Categorical).
        \item Ưu tiên 3 -- Màu sắc (Hue Color): Việc dùng Hue để tô màu cho các cột là một dạng mã hóa lặp (Redundant Encoding). Tuy không bắt buộc do đã có PosX phân tách, nhưng nó đóng vai trò liên kết thị giác (nhấn mạnh Xanh = Cá nhân, Cam = Chuyên nghiệp) xuyên suốt từ biểu đồ 9.1 sang.
    \end{itemize}
\end{itemize}

\textbf{2. Nguyên lý hiệu quả:}

\begin{itemize}
    \item Độ chính xác (Accuracy): Trục PosY biểu diễn thông qua chiều dài (Length) bám sát gốc 0. Theo định luật Stevens, Length có $n = 1.0$, do đó mắt người cảm nhận cực kỳ chính xác. Độ lỗi Log\_error rất thấp.
    \item Khả năng phân biệt (Discriminability): 2 cột đứng cạnh nhau với 2 màu tương phản mạnh (Xanh/Cam) giúp mắt người đọc ngay lập tức nhận ra bên nào cao hơn mà không cần suy nghĩ.
    \item Khả năng tách biệt (Separability): Sử dụng Facet (Juxtapose) giải quyết việc tách biệt 2 hệ đo lường khác nhau (Rating thang 5, IsBooked là tỷ lệ \%) $\Rightarrow$ không xảy ra hiện tượng nhiễu trục.
\end{itemize}

\subsubsection*{C. Phân tích biểu đồ (Insight)}

\begin{itemize}
    \item Dù yếu thế về số lượng nguồn cung, chủ nhà ``Cá nhân'' (màu xanh) lại mang đến trải nghiệm tốt hơn so với nhóm chuyên nghiệp.
    \item Hệ quả là tỷ lệ lấp đầy của chủ nhà ``Cá nhân'' cũng cao hơn nhóm chuyên nghiệp. Điều này chứng minh khách hàng đề cao tính cá nhân hóa, sự thân thiện của các cá thể hơn là dịch vụ rập khuôn của các tổ chức.
\end{itemize}
